一般超胞由一個整數矩陣 $S$ 定義。矩陣的每一行說明一條新晶格向量由原來的三條向量怎樣組合。例如，第一行為 $(2,1,0)$，就表示 $\mathbf a'=2\mathbf a+\mathbf b$。在本書的行向量約定下，晶格、體積和位點數滿足
\begin{equation}
 A'=SA,\qquad V'=|\det S|V,\qquad N'=|\det S|N.
 \label{eq:supercell-count-volume}
\end{equation}
$S$ 中的數必須是整數，且行列式不能為零。$|\det S|$ 是新舊晶胞的體積比，也是獨立位點數的倍數；$N$ 始終數模型位點，不數邊界上額外繪出的原子映像。對於三個重複數，$S=\operatorname{diag}(n_a,n_b,n_c)$，式\eqref{eq:supercell-count-volume}就回到三個重複數相乘的規則。前面的矽超胞例子取 $S=2I$，所以 $|\det S|=8$，8個位點變為64個。

下面用一個非對角矩陣作代數檢查：
\[
 S=\begin{bmatrix}2&1&0\\0&1&0\\0&0&1\end{bmatrix},\qquad
 \begin{cases}
 \mathbf a'=2\mathbf a+\mathbf b,\\
 \mathbf b'=\mathbf b,\quad \mathbf c'=\mathbf c.
 \end{cases}
\]
其行列式為2。若原模型含8個位點，新模型應含16個位點，體積是原來的兩倍；新外框可以傾斜，週期材料本身的區域性排列仍來自原結構的複製。這與透過應變拉伸原子間距有區別：超胞選擇更大的重複單元，不把每條原子間連線都拉長兩倍。

在\ui{超胞倍率或變換矩陣}中，可以填寫三個重複整數，也可以填寫完整的 $3\times3$ 矩陣。使用一般矩陣前先寫出三條新基矢，再檢查行列式、預覽外框和應用後的位點數，便不容易把“重新定向”和“增加了多少原子”混為一談。輸入\literalinput{1 1 1}時，矩陣為 $I$，模型不擴大；本工具仍會把等價座標歸入單位晶胞，適合將落在胞外的代表位置整理到同一胞中檢視。
